\documentclass[12pt,a4paper]{article}

\usepackage[margin=2.5cm]{geometry}
\usepackage{amsmath}
\usepackage{booktabs}
\usepackage{hyperref}
\usepackage{graphicx}
\usepackage{xcolor}
\usepackage{array}
\usepackage{caption}
\usepackage{subcaption}
\usepackage{parskip}

\hypersetup{
    colorlinks=true,
    linkcolor=blue,
    urlcolor=blue,
    citecolor=blue
}

\definecolor{passgreen}{RGB}{0,150,0}
\definecolor{failred}{RGB}{200,0,0}
\definecolor{warnorange}{RGB}{200,120,0}

\title{pyCSEP Catalog-Based Evaluation:\\
       NZ ETAS Forecast --- 2016 Kaik\={o}ura Sequence}
\author{NZ ETAS Pipeline}
\date{\today}

% ────────────────────────────────────────────────────────────────────────────
\begin{document}
\maketitle
\tableofcontents
\newpage

% ────────────────────────────────────────────────────────────────────────────
\section{Overview and Data Provenance}

This report presents a full pyCSEP (v0.8.0) catalog-based evaluation of the
NZ ETAS forecast for the 2016 M7.82 Kaik\={o}ura earthquake sequence
(GeoNet event \texttt{2016p858000}, mainshock UTC origin time
\texttt{2016-11-13T11:02:56Z}).

\subsection{Two-phase protocol: calibration then forecast}

The pipeline operates in two strictly separate phases, configured in
\texttt{etas\_config.json}:

\begin{description}
    \item[\textbf{Data / calibration window} (Days 0--7 post-mainshock):]
        GeoNet aftershocks from
        \texttt{2016-11-13T11:02:56Z} to \texttt{2016-11-20T11:02:56Z}
        within a 200\,km epicentral radius are fetched
        (\texttt{dataWindow.minDays = 0}, \texttt{maxDays = 7}).
        These observations are used \emph{only} to fit the ETAS
        parameters via maximum-likelihood grid search;
        they are never compared against the forecast.

    \item[\textbf{Forecast window} (Days 7--14 post-mainshock):]
        The fitted model is run forward from
        \texttt{2016-11-20T11:02:56Z} to \texttt{2016-11-27T11:02:56Z}
        (\texttt{forecastWindow.minDays = 7}, \texttt{maxDays = 14}).
        All 1{,}000 Monte Carlo simulations start at $t = 7$~days; the
        simulated catalogs confirm this: \texttt{sim\_0008.txt} line 4
        begins at $t = 7.001$~days, not $t = 0$.
        The GeoNet observed catalog over the same absolute window is
        fetched independently and used only for evaluation.
\end{description}

\subsection{ETAS configuration (\texttt{etas\_config.json})}

\begin{description}
    \item[\textbf{Generic NZ priors}:]
        $\bar{a} = -2.423\ (\sigma = 0.395)$,
        $\bar{p} = 1.08\ (\sigma = 0.20)$,
        $\bar{c} = 0.01\ (\sigma_{\log c} = 0.70)$,
        $\alpha = 1.0$, $b = 1.0$, $M_\text{ref} = 4.5$.
        These are the standard NZ generic ETAS parameters.

    \item[\textbf{MLE grid search}:]
        $a_\text{ms}$: 21 steps in $[-4,\ 1]$;
        $a$: 21 steps in $[-3.5,\ -0.5]$;
        $p$: 16 steps in $[0.5,\ 2.0]$;
        $c$: 16 log-spaced steps in $[10^{-4},\ 10^{-1}]$.
        Total grid: $21 \times 21 \times 16 \times 16 = 112{,}896$ points.
        \texttt{fitMSProductivity = true} fits the mainshock productivity
        parameter $a_\text{ms}$ separately from the sequence $a$.

    \item[\textbf{Monte Carlo simulation}:]
        $N_\text{sims} = 1{,}000$, maximum 10 generations,
        $M_\text{max} = 9.5$, \texttt{seed = 12345}.

    \item[\textbf{Spatial kernel}:]
        $0.05°$ grid spacing, length scale = 5.0, stress drop = 3.0~MPa,
        \texttt{fitType = "aftershocks"} (inferred line-source geometry
        from aftershock distribution). Each simulated event is assigned a
        location sampled from the isotropic ETAS spatial kernel
        $S(r) \propto (d^2 + r^2)^{-1}$ using exact analytical CDF inversion.
\end{description}

\textbf{Script}: \texttt{scripts/python/visualize\_nz\_etas\_with\_pycsep.py}

\subsection{What the evaluation script does}

\textbf{Steps, in order}:
\begin{enumerate}
    \item Reads the ETAS summary file (\texttt{nz\_etas\_simulations.txt}) and
          config (\texttt{etas\_config.json}) to extract the event ID, forecast
          window (Days 7.0--14.0), and $M_c = 3.0$.
    \item Fetches the mainshock origin time from GeoNet FDSN and caches it.
    \item Loads all \texttt{sim\_*.txt} files from \texttt{simulated\_catalogs/}.
          Filters each catalog to $M \geq M_c = 3.0$, then maps events into
          the official
\newline \textbf{pyCSEP NZ testing region}
          (\texttt{regions.nz\_csep\_region()}) using \texttt{filter\_spatial}.
    \item Queries GeoNet for the observed catalog over the same window and
          spatial region, caches it.
    \item Builds a \texttt{CatalogForecast} from the filtered simulations and
          a \texttt{GriddedForecast} of expected rates via
          \texttt{forecast.get\_expected\_rates()}.
    \item Runs six catalog-based consistency tests (N, M, S, PL, resampled M,
          MLL) from \texttt{csep.core.catalog\_evaluations}.
    \item Runs a rolling 1-day N-test and a calibration test across the seven
          daily sub-windows.
    \item Produces all figures and saves JSON results to
          \texttt{build/pycsep/evaluation\_json/}.
\end{enumerate}

\textbf{Key run statistics} (from \texttt{nz\_etas\_pycsep\_summary.md}):
\begin{itemize}
    \item 1000 simulated catalogs loaded.
    \item 3{,}551{,}889 raw event rows read; 360{,}343 kept after $M \geq 3.0$
          filter; 348{,}320 kept after NZ pyCSEP region filter.
    \item Ensemble event-count \textbf{median: 334} events.
    \item Median indirect ($\text{Generation} \geq 2$) share: \textbf{32.6\%}.
    \item Observed GeoNet catalog: \textbf{$N_\text{obs} = 323$} events
          ($M_w \geq 3.0$, max $M_w = 5.81$,
          \texttt{2016-11-20T11:02:56Z} to \texttt{2016-11-27T11:02:56Z} ---
          the forecast window only, not the 0--7~day calibration period).
    \item Representative catalog: \texttt{sim\_0102} (334 events --- closest
          to the ensemble median, selected by
          \texttt{representative\_catalog()} which minimises
          $|N_\text{sim} - \text{median}|$).
\end{itemize}

\textbf{Spatial region note}: The pyCSEP NZ region
(\texttt{regions.nz\_csep\_region()}) is a fixed official testing grid.
3{,}002 of 6{,}343 cells in the region had zero ETAS spatial rate --- the model
assigns non-zero probability to roughly half the region. Sixteen observed
events fell in zero-rate cells; 51 observed events fell in zero
space--magnitude bins. This triggered \texttt{status=undersampled} for the S
and PL tests (the score was recomputed after excluding those events).

% ────────────────────────────────────────────────────────────────────────────
\section{Key Statistics}

\subsection{Delta statistics ($\delta_1$, $\delta_2$) --- N-test only}

\begin{align}
    \delta_1 &= 1 - F(N_\text{obs} - 1 \mid N_F)
               \quad \text{probability of observing \emph{at least} $N_\text{obs}$ events} \\[4pt]
    \delta_2 &= F(N_\text{obs} \mid N_F)
               \quad \text{probability of observing \emph{at most} $N_\text{obs}$ events}
\end{align}

pyCSEP uses the \emph{empirical} distribution of synthetic catalog sizes (not
a Poisson CDF) for catalog-based forecasts, so $F$ here is the empirical CDF
of the 1000 simulation counts. Both $\delta_1$ and $\delta_2$ must exceed
0.025 to pass at 95\% confidence.

\subsection{Quantile scores for all other tests}

Each test returns a two-element quantile tuple $(q_\text{lo}, q_\text{hi})$
where $q_\text{lo} = P(\text{sim stat} \leq \omega_\text{obs})$ and
$q_\text{hi} = 1 - q_\text{lo}$. For one-sided tests the relevant tail is
$q_\text{lo}$: a value below 0.05 indicates the observation is less likely
than 95\% of simulations --- a borderline or failing result.

% ────────────────────────────────────────────────────────────────────────────
\section{Figure 1 --- Evaluation Distributions (Six-Panel)}
\label{sec:eval_dist}

\begin{figure}[h!]
    \centering
    \includegraphics[width=\textwidth]{OAF_figures/nz_etas_pycsep_evaluation_distributions.png}
    \caption{pyCSEP catalog-based evaluation distributions (1000 simulations,
             $M_w \geq 3.0$, Days 7.0--14.0). Each panel plots the
             distribution of a test statistic across the 1000 simulated
             catalogs (blue histogram). The dashed vertical line is the
             observed statistic $\omega$. Red bars are the 2.5\% tails.
             The figure is produced by \texttt{plots.plot\_test\_distribution}
             called once per result, then assembled into a $2 \times 3$ grid
             by the script. Test results are drawn from
             \texttt{csep.core.catalog\_evaluations}.}
    \label{fig:eval}
\end{figure}

\subsection{Number Test (N-test) --- top left}

\textbf{What the script computes}:
\newline
\texttt{catalog\_evaluations.number\_test(forecast, observed\_catalog)}.
The statistic is simply $N_\text{obs} = 323$. The test distribution is the
empirical distribution of the 1000 simulated catalog sizes (post NZ-region
filter), ranging from 230 to 616 with median 334.

\textbf{Exact values from \texttt{catalog\_number\_test.json}}:
\[
    \omega = 323, \quad
    \delta_1 = 0.594, \quad \delta_2 = 0.411, \quad \text{status: normal}
\]
Both exceed 0.025. The observation sits below the simulation median (334),
meaning GeoNet recorded slightly fewer events than the typical simulation.
\textbf{Status}: \textcolor{passgreen}{\textbf{Pass}}

\subsection{Magnitude Test (M-test) --- top centre}

\textbf{What the script computes}:
\newline
\texttt{catalog\_evaluations.magnitude\_test(forecast, observed\_catalog)}.
Sums all simulated catalogs over spatial cells to get a magnitude-only
union distribution, then computes the joint log-likelihood of the observed
magnitude sequence given that union.

\textbf{Exact values}:
\[
    \omega = 0.5501, \quad
    q_\text{lo} = 0.736, \quad q_\text{hi} = 0.264, \quad \text{status: normal}
\]
73.6\% of simulations scored \emph{below} (worse than) the observation.
The $b = 1$ GR magnitude distribution is well-calibrated.
\textbf{Status}: \textcolor{passgreen}{\textbf{Pass}}

\subsection{Spatial Test (S-test) --- top right}

\textbf{What the script computes}:
\newline
\texttt{catalog\_evaluations.spatial\_test(forecast, observed\_catalog)}.
Sums over magnitude bins to get a spatial-only distribution. Computes the
mean log-likelihood of observed event locations under the forecast's spatial
density. Because 16 observed events fell in zero-rate cells, pyCSEP removed
them and recomputed (\texttt{status=undersampled}).

\textbf{Exact values}:
\[
    \omega = -5.739, \quad
    q_\text{lo} = 0.963, \quad q_\text{hi} = 0.037, \quad \text{status: undersampled}
\]
The relevant tail for this one-sided test is $q_\text{hi} = 0.037 < 0.05$.
Only 3.7\% of simulations scored \emph{worse} spatially than the observation.
The observed events cluster in locations the model assigns lower probability
than expected --- the isotropic kernel misses part of the complex multi-fault
Kaik\={o}ura rupture zone.
\textbf{Status}: \textcolor{warnorange}{\textbf{Marginal fail ($q = 0.037$)}}

\subsection{Pseudolikelihood Test (PL-test) --- bottom left}

\textbf{What the script computes}:
\texttt{catalog\_evaluations.pseudolikelihood\_test(forecast, observed\_catalog)}.
Evaluates the joint space--magnitude pseudo-log-likelihood:
\[
    \hat{L}_\text{obs} = \sum_i \log \hat{\lambda}_s(k_i) - \bar{N}
\]
Also \texttt{status=undersampled} for the same reason as S-test.

\textbf{Exact values}:
\[
    \omega = -313.41, \quad
    q_\text{lo} = 0.969, \quad q_\text{hi} = 0.031, \quad \text{status: undersampled}
\]
$q_\text{hi} = 0.031 < 0.05$ --- also marginal. The joint
space--magnitude likelihood is slightly low because of the same
spatial mismatch identified by the S-test.
\textbf{Status}: \textcolor{warnorange}{\textbf{Marginal fail ($q = 0.031$)}}

\textbf{Note}: pyCSEP recommends the PL-test as the primary joint test for
catalog-based forecasts. Both it and the S-test share the same root cause
(the undersampled spatial cells).

\subsection{Resampled Magnitude Test --- bottom centre}

\textbf{What the script computes}:
\texttt{catalog\_evaluations.resampled\_magnitude\_test(forecast,
observed\_catalog, seed=12345)}.
Draws exactly $N_\text{obs} = 323$ events from the union of all synthetic
catalogs (with fixed \texttt{seed=12345} for reproducibility) to build each
test distribution entry.

\textbf{Exact values}:
\[
    \omega = 0.5501, \quad
    q_\text{lo} = 0.821, \quad q_\text{hi} = 0.179, \quad \text{status: normal}
\]
The wider spread relative to the plain M-test reflects the additional
sampling variance from resampling.
\textbf{Status}: \textcolor{passgreen}{\textbf{Pass}}

\subsection{MLL Magnitude Test --- bottom right}

\textbf{What the script computes}:
\texttt{catalog\_evaluations.MLL\_magnitude\_test(forecast, observed\_catalog,
seed=12345)}.
Applies a modified multinomial log-likelihood ratio statistic to the
magnitude distribution.

\textbf{Exact values}:
\[
    \omega = 163.598, \quad
    q_\text{lo} = 0.356, \quad q_\text{hi} = 0.644, \quad \text{status: normal}
\]
The extremely narrow distribution (values $\approx 160$--$166$) reflects
how tightly the GR law constrains magnitude statistics across all simulations.
\textbf{Status}: \textcolor{passgreen}{\textbf{Pass}}

% ────────────────────────────────────────────────────────────────────────────
\section{Figure 2 --- Rolling Number-Test Consistency}
\label{sec:rolling}

\begin{figure}[h!]
    \centering
    \includegraphics[width=0.85\textwidth]{OAF_figures/nz_etas_pycsep_rolling_number_consistency.png}
    \caption{Rolling N-test consistency across the seven consecutive 1-day
             sub-windows of the forecast period (Days 7.0--8.0, 8.0--9.0,
             \ldots, 13.0--14.0 post-mainshock; i.e.\ 20--27 November 2016
             UTC). All times are relative to the mainshock origin
             \texttt{2016-11-13T11:02:56Z}. Produced by
             \texttt{run\_rolling\_number\_tests()} which builds a
             \texttt{CatalogForecast} and an observed sub-catalog for each
             window, calls \texttt{catalog\_evaluations.number\_test}, and
             plots the result with \texttt{plots.plot\_number\_test}.
             Each row is one day. The horizontal line and shaded band are the
             95\% CI of simulated daily counts. The dot is the simulation mean.
             Green squares = observed count inside CI (pass);
             red dots = outside CI (fail). All seven days pass.
             The rolling calibration KS statistic is 0.160 with $p = 0.980$,
             confirming the seven sub-window tests are jointly well-calibrated.}
    \label{fig:rolling}
\end{figure}

\textbf{What the script computes}: For each of the 7 daily sub-windows, the
script slices both the forecast catalogs and the observed catalog to that
window using 
\newline
\texttt{filter\_catalog\_absolute\_window}, then calls
\texttt{catalog\_evaluations.number\_test}. A rolling calibration test 
(\texttt{catalog\_evaluations.calibration\_test}) is then run across all 7
results to test whether the sub-window p-values are uniformly distributed
(KS statistic = 0.160, $p = 0.980$ --- consistent with uniform, meaning
no systematic temporal bias).

All seven daily windows show green squares. The Omori temporal decay is
well-calibrated not just in aggregate but day by day. The broad CI
($\sim$20--100 events/day) reflects ETAS branching overdispersion.

% ────────────────────────────────────────────────────────────────────────────
\section{Figure 3 --- Spatial Expected Rate Map}
\label{sec:spatial_map}

\begin{figure}[h!]
    \centering
    \includegraphics[width=0.85\textwidth]{OAF_figures/nz_etas_pycsep_expected_rates.png}
    \caption{Expected $M \geq 3.0$ rate map averaged over 1000 simulations,
             expressed as $\log_{10}$ events per NZ pyCSEP grid cell per
             0.019 years (the 7-day window). Produced by
             \texttt{forecast.get\_expected\_rates()} which averages the
             gridded counts across all 1000 \texttt{CSEPCatalog} objects
             filtered into the NZ region, then calls \texttt{expected\_rates.plot()}.
             The NZ pyCSEP region covers 34°S--48°S, 166°E--179°E on a
             $0.1°\times 0.1°$ grid (6{,}343 cells). Colour scale: purple ($-3$,
             $\sim$0.001 events/cell) to yellow ($+1$, $\sim$10 events/cell).
             The dominant hot spot near 42°S, 173°E is the Kaik\={o}ura aftershock
             zone; its NE--SW elongation reflects the inferred line-source geometry
             (\texttt{fitType = "aftershocks"}).}
    \label{fig:spatial_map}
\end{figure}

\textbf{What the script computes}: \texttt{forecast.get\_expected\_rates(verbose=False)}
accumulates all $M \geq M_c$ events from the 1000 simulated catalogs
(already filtered into the NZ region) into a \texttt{GriddedForecast}.
Each cell value is the mean number of events across the ensemble. Cells
where no simulation placed any event remain at zero --- this is what
produces the 3{,}002 zero-rate cells noted in Section~1, and is the
root cause of the S-test undersampling.

The map directly shows why the S-test is marginal: the smooth, roughly
elliptical high-rate zone covers only part of the actual aftershock
distribution. The Kaik\={o}ura sequence ruptured $\geq 12$ separate fault
segments; the isotropic kernel cannot capture that irregular geometry.

% ────────────────────────────────────────────────────────────────────────────
\section{Figure 4 --- Magnitude Histogram}
\label{sec:mag_hist}

\begin{figure}[h!]
    \centering
    \includegraphics[width=0.85\textwidth]{OAF_figures/nz_etas_pycsep_histogram.png}
    \caption{Magnitude histogram comparing the GeoNet observed catalog
             (blue dots, $N_\text{obs} = 323$) against the ETAS ensemble
             forecast (red dots = mean; red bars = 95\% CI). Produced by
             \texttt{plots.plot\_magnitude\_histogram(forecast,
             observed\_catalog, log\_scale=False, normalize=False)}.
             The x-axis covers $M_w$ 3.0--9.5 in 0.1 unit bins; y-axis is
             raw event count. The GR $b = 1$ slope is clearly reproduced.
             All observed points fall within the 95\% CI.}
    \label{fig:mag_hist}
\end{figure}

\textbf{What the script computes}: \texttt{plots.plot\_magnitude\_histogram}
bins both the union of simulated catalogs and the observed catalog in 0.1
magnitude units. The CI bars come from the per-simulation counts at each bin
(5th and 95th percentiles across 1000 simulations). The ``comparison
catalog'' passed to this function is the live GeoNet observed catalog when
available (which it is here).

Key observations:
\begin{itemize}
    \item At $M3.0$--$3.1$, $\approx 75$ observed events vs.\ $\approx 85$
          ensemble mean --- slight under-observation consistent with
          short-term catalog incompleteness in the days immediately
          post-mainshock.
    \item The GR slope ($b = 1$) is reproduced at all magnitudes.
    \item The M-test quantile ($q_\text{lo} = 0.736$) confirms the
          magnitude distribution is better fit than 73.6\% of simulations.
\end{itemize}

% ────────────────────────────────────────────────────────────────────────────
\section{Figure 5 --- Cumulative Count Time Series}
\label{sec:cumulative}

\begin{figure}[h!]
    \centering
    \includegraphics[width=0.85\textwidth]{OAF_figures/nz_etas_pycsep_cumulative.png}
    \caption{Cumulative $M \geq 3.0$ event counts over the 7-day forecast
             window (Days 7--14 post-mainshock; absolute dates
             2016-11-20 to 2016-11-27 UTC). The x-axis is 0--7 days
             \emph{relative to forecast start} (i.e.\ day~7 absolute),
             not relative to the mainshock. Produced by
             \texttt{plots.plot\_cumulative\_events\_versus\_time(forecast,
             observed\_catalog, time\_axis="days")}.
             The black line is the GeoNet observed catalog
             ($N_\text{obs} = 323$ at the end of the window). The dark
             blue band is the ensemble 25th--75th percentile (IQR); the
             light blue band is the 5th--95th percentile. The ensemble
             median is 334. The observation tracks near the 40th--60th
             percentile throughout the window, consistent with
             $\delta_1 = 0.594$, $\delta_2 = 0.411$.}
    \label{fig:cumulative}
\end{figure}

\textbf{What the script computes}: The cumulative plot uses the live observed
catalog (not the representative simulation) because \texttt{observed\_catalog
is not None}. The function evaluates the cumulative count of each of the 1000
simulations and the observation at a fine time grid, then plots quantile bands.
The Omori decay is visible as a decreasing slope through the window.

Key observations:
\begin{itemize}
    \item The observed line stays inside the IQR for the full window.
    \item By day 7, the observation (323) is below the median (334), placing
          it at roughly the 41st percentile --- consistent with $\delta_2 = 0.411$.
    \item The upper 95th-percentile tail extends far above the median
          (characteristic ETAS branching overdispersion from rare
          supercritical cascades). The 32.6\% indirect ($\text{Gen} \geq 2$)
          share confirms multi-generation triggering is active.
    \item No visible sudden jump in the observed curve, consistent with the
          moderate observed maximum $M_w = 5.81$.
\end{itemize}

% ────────────────────────────────────────────────────────────────────────────
\section{Summary of All Tests}

\begin{table}[h!]
\centering
\caption{pyCSEP evaluation results for the 2016 Kaik\={o}ura ETAS forecast
         (1000 simulations, $M_w \geq 3.0$, Days 7.0--14.0, pyCSEP v0.8.0).
         All values from \texttt{nz\_etas\_pycsep\_summary.md}.}
\begin{tabular}{p{3.8cm} p{2.8cm} p{2.5cm} p{2.2cm} p{2cm}}
\toprule
\textbf{Test} & \textbf{Statistic $\omega$} & \textbf{Quantile} & \textbf{Threshold} & \textbf{Status} \\
\midrule
N-test (full)          & 323               & $(0.594,\ 0.411)$ & both $>0.025$  & \textcolor{passgreen}{\textbf{Pass}}      \\
N-test (rolling, 7d)   & KS $= 0.160$      & $p = 0.980$       & ---            & \textcolor{passgreen}{\textbf{Pass}}      \\
M-test                 & $0.5501$          & $(0.736,\ 0.264)$ & $q_\text{hi}>0.05$ & \textcolor{passgreen}{\textbf{Pass}}  \\
S-test$^\dagger$       & $-5.739$          & $(0.963,\ 0.037)$ & $q_\text{hi}>0.05$ & \textcolor{warnorange}{\textbf{Marginal}} \\
PL-test$^\dagger$      & $-313.41$         & $(0.969,\ 0.031)$ & $q_\text{hi}>0.05$ & \textcolor{warnorange}{\textbf{Marginal}} \\
Resampled M-test       & $0.5501$          & $(0.821,\ 0.179)$ & $q_\text{hi}>0.05$ & \textcolor{passgreen}{\textbf{Pass}}  \\
MLL magnitude          & $163.598$         & $(0.356,\ 0.644)$ & supplementary  & \textcolor{passgreen}{\textbf{Pass}}      \\
\bottomrule
\end{tabular}

\smallskip
\noindent $^\dagger$ \texttt{status=undersampled}: 16 observed events fell in
zero-rate spatial cells and were excluded before recomputing the score.
\end{table}

\subsection{Overall interpretation}

The S-test and PL-test are both marginal ($q = 0.037$ and $q = 0.031$
respectively), sharing the same root cause: the isotropic spatial kernel
assigns zero rate to cells where 16 of 323 observed events (5\%) actually
occurred. Both tests are flagged \texttt{undersampled} by pyCSEP, which
recomputes the score after excluding those events --- the residual marginal
score reflects the remaining spatial mismatch in cells that do have ETAS rate.

The temporal (N-test) and magnitude (M, resampled-M, MLL) components all
pass cleanly. The rolling calibration confirms no systematic temporal bias
($p = 0.980$).

The primary fix is the spatial kernel: switching from
\texttt{fitType = "aftershocks"} (inferred line source) to
\texttt{fitType = "shakemap"} (geodetically constrained fault geometry) would
distribute aftershock probability along all $\geq 12$ fault segments, filling
in the zero-rate cells and likely bringing both S and PL tests into the
passing range.

% ────────────────────────────────────────────────────────────────────────────
\section*{References}
\begin{itemize}
    \item Savran, W.\ et al.\ (2022). pyCSEP: A Python Toolkit for Earthquake
          Forecast Developers. \textit{Seismological Research Letters}, 93(5).
          \url{https://doi.org/10.1785/0220210290}
    \item Schorlemmer, D.\ et al.\ (2007). Earthquake likelihood model testing.
          \textit{Seismological Research Letters}, 78(1), 17--29.
    \item Zechar, J.D.\ et al.\ (2010). Likelihood-based tests for evaluating
          space--rate--magnitude earthquake forecasts.
          \textit{Bulletin of the Seismological Society of America}, 100(3),
          1184--1195.
    \item pyCSEP documentation: \url{https://docs.cseptesting.org}
\end{itemize}

\end{document}
